\chapter{Đọc thêm 2: Sử dụng tư duy phản biện để nhận diện tin giả (Machete \& Turpin)}
\label{ch:M5L4DocThem2}

\section{Thông tin nguồn}

\begin{itemize}
  \item \textbf{Nguồn:} Machete, P. \& Turpin, M. ``The Use of Critical Thinking to Identify Fake News: A Systematic Literature Review.'' Springer, 2020.
  \item \textbf{Phần cần đọc:} Toàn bộ bài báo
  \item \textbf{Ghi chú:} bản chuyển ngữ tiếng Việt súc tích, bám cấu trúc nguồn (giữ nguyên tiêu đề, công thức, số liệu, bảng, chú thích và danh mục tài liệu tham khảo).
\end{itemize}

Việc sử dụng tư duy phản biện để nhận diện tin giả: Một tổng quan tài liệu có hệ thống

Paul Machete và Marita Turpin Khoa Tin học, Đại học Pretoria, Pretoria 0001, Nam Phi marita.turpin@up.ac.za

\textbf{Tóm tắt.} Với lượng tin tức khổng lồ được đăng tải trực tuyến hiện nay, khả năng đánh giá độ tin cậy của tin tức trực tuyến đã trở nên thiết yếu. Dù đã có nhiều nghiên cứu về tin giả và các công cụ phát hiện nó, số công trình tập trung vào việc dùng năng lực thông tin (information literacy) để giúp con người tiếp cận thông tin và tin tức trực tuyến một cách phản biện vẫn còn hạn chế. Tư duy phản biện (critical thinking), với tư cách một dạng năng lực thông tin, cung cấp phương thức tiếp cận nội dung trực tuyến có phê phán, chẳng hạn tìm bằng chứng cho các luận điểm và đánh giá tính hợp lý của lập luận. Mục đích của nghiên cứu này là khảo sát hiện trạng tri thức về việc dùng tư duy phản biện để nhận diện tin giả. Nhóm tác giả thực hiện một tổng quan tài liệu có hệ thống (SLR) nhằm xác định các nghiên cứu trước đây về đánh giá độ tin cậy của tin tức, đặc biệt là việc sử dụng tư duy phản biện để đánh giá tin tức trực tuyến. Qua quá trình sàng lọc của SLR, 22 nghiên cứu phù hợp được xác định. Mặc dù một số nghiên cứu có đề cập năng lực thông tin, chỉ ba nghiên cứu bàn trực tiếp về tư duy phản biện như một phương tiện nhận diện tin giả. Các nghiên cứu này coi tư duy phản biện là kỹ năng thiết yếu để nhận diện tin giả và khuyến nghị đưa giáo dục năng lực thông tin vào các cơ sở học thuật, đặc biệt nhằm khuyến khích tư duy phản biện.

\textbf{Từ khóa:} Tư duy phản biện, Tin giả, Năng lực thông tin, Tổng quan tài liệu có hệ thống

\section{1 Giới thiệu}

Thời đại thông tin đã làm tăng đáng kể các nguồn thông tin sẵn có, song hành với sự gia tăng chưa từng có của khả năng truy cập và kết nối internet cùng mức độ phổ biến của các thiết bị công nghệ {[}1{]}. Con người không còn chỉ dựa vào truyền hình và báo in để theo dõi tin tức mà ngày càng dùng mạng xã hội và ứng dụng tin tức. Sự đa dạng của các nguồn thông tin hiện nay đã góp phần làm lan truyền những "sự thật thay thế" (alternative facts) {[}1{]}. Với hơn 1,8 tỷ người dùng hoạt động mỗi tháng vào năm 2016 {[}2{]}, Facebook chiếm 20\% tổng lưu lượng truy cập vào các trang web đáng tin cậy và tới 50\% tổng lưu lượng truy cập vào các trang tin giả {[}3{]}. Twitter đứng thứ hai sau Facebook, với hơn 400 triệu người dùng hoạt động mỗi tháng {[}2{]}. Các bài đăng trên mạng xã hội như Facebook và Twitter lan truyền nhanh vì chúng tìm cách thu hút sự chú ý của người đọc càng nhanh càng tốt, với rất ít thông tin thực chất, nhờ đó tạo mảnh đất màu mỡ cho việc phát tán tin giả {[}4{]}.

Mạng xã hội là cách tiện lợi để tiếp cận tin tức và giữ liên lạc với bạn bè, gia đình, nhưng không dễ phân biệt tin thật với tin giả trên đó {[}5{]}. Mạng xã hội tiếp tục góp phần làm gia tăng việc phân phối thông tin do người dùng tạo ra, bao gồm tin lừa bịa, các tuyên bố sai sự thật, tin bịa đặt và thuyết âm mưu, với nguồn chính là các nền tảng như Facebook và Twitter {[}6{]}. Điều này có nghĩa là bất kỳ ai sở hữu thiết bị kết nối internet đều có thể là người tiêu thụ hoặc người phát tán tin giả. Các nền tảng mạng xã hội và công cụ tìm kiếm tuy không khuyến khích người dùng tin vào thông tin lưu hành, nhưng vẫn là đồng lõa trong xu hướng người dùng tin vào thông tin gặp phải trên các nền tảng này mà không xác minh tính xác thực {[}6{]}. Sự lan truyền của tin giả có thể gây nhiều thiệt hại cho đối tượng bị nhắm tới, từ tổn hại danh tiếng của một cá nhân đến ảnh hưởng tới giá trị cảm nhận của một công ty {[}7{]}.

Mục đích của nghiên cứu này là khảo sát việc sử dụng các phương pháp tư duy phản biện để phát hiện các tin không đúng sự thật, hoặc giúp hình thành thái độ phản biện đối với tin tức trực tuyến. Công trình được thực hiện bằng tổng quan tài liệu có hệ thống (SLR). Bài báo được bố cục như sau. Phần tiếp theo cung cấp thông tin nền về tin giả, tầm quan trọng của nó trong đời sống hằng ngày của người dùng mạng xã hội, và cách năng lực thông tin cùng tư duy phản biện có thể được dùng để nhận diện tin giả. Sau đó, phương pháp nghiên cứu SLR được trình bày. Tiếp theo, các kết quả của tổng quan được báo cáo, trước hết theo thống kê mô tả và sau đó theo phân tích chủ đề của các nghiên cứu đã xác định. Bài báo kết thúc bằng phần Kết luận và các khuyến nghị.

\section{2 Bối cảnh: Tin giả, năng lực thông tin và tư duy phản biện}

Phần này bàn về lịch sử của tin giả, tin giả như chúng ta biết ngày nay, và vai trò của năng lực thông tin trong việc hỗ trợ nhận diện tin giả. Phần này cũng nêu ngắn gọn định nghĩa về tư duy phản biện.

\subsection{2.1 Lịch sử của tin giả}

Dù tin giả (fake news) mới được chú ý nhiều gần đây, thuật ngữ này đã được giới học thuật sử dụng từ lâu {[}4{]}. Tin giả bắt nguồn từ truyền thống báo chí giật gân (yellow journalism) của thập niên 1890, vốn dựa vào các yếu tố quen thuộc của chủ nghĩa giật gân: tin tội phạm, bê bối và đồn đại, ly hôn và tình dục, nhấn mạnh việc đưa tin thảm họa, thể thao giật gân, và có thể cả tin châm biếm {[}5{]}. Sự xuất hiện của tin tức trực tuyến đầu những năm 2000 làm dấy lên lo ngại, trong đó có việc những người cùng quan điểm có thể hình thành các "buồng vang" (echo chambers) để lọc bỏ những ý kiến khác {[}2{]}. Đó là hệ quả của việc truyền thông chuyển từ mô hình do báo in của các nhà báo chính danh, đáng tin cậy chi phối sang mô hình trong đó nhiều người tin vào tin tức trực tuyến từ nguồn thiếu tin cậy {[}5{]}. Về sau, thuật ngữ này còn dùng để chỉ các "chương trình tin tức châm biếm", "chương trình tin tức nhại" hay "chương trình hài tin giả", tức các chương trình truyền hình hoặc một phần của chương trình dành cho châm biếm chính trị {[}4{]}. Có thể kể đến The Daily Show (nay do Trevor Noah dẫn), mục "The Weekend Update" của Saturday Night Live, cùng các chương trình tương tự như Last Week Tonight with John Oliver và The Colbert Report của Stephen Colbert {[}4{]}. Các bản tin trong những chương trình này bị gọi là "giả" không phải vì nội dung, mà vì chúng nhại lại tin tức của các đài truyền hình bằng giọng mỉa mai và dùng hài kịch làm công cụ tiếp cận các vấn đề công chúng có thật {[}4{]}. Cụm từ "Fake News" càng trở nên nổi bật trong cuộc bầu cử tổng thống Mỹ năm 2016, khi các thành viên của các đảng đối lập đăng tiêu đề tin sai sự thật nhằm tác động đến quyết định của cử tri {[}6{]}.

\subsection{2.2 Tin giả ngày nay}

Ngày nay thuật ngữ tin giả mang nghĩa đen rõ hơn {[}4{]}. Từ điển Macquarie chọn "fake news" là từ của năm 2016 {[}8{]}; theo từ điển này, đó là từ phản ánh một bước tiến đáng chú ý trong việc tạo ra nội dung lừa dối, đồng thời cho phép mỗi người tin vào điều mình cho là hợp lý. Có nhiều định nghĩa cho cụm từ này, nhưng một mô tả ngắn gọn là của Paskin {[}4{]}: những bài báo xuất phát từ mạng xã hội hoặc các nền tảng chính thống (trực tuyến hay ngoại tuyến), không dựa trên sự thật nhưng được trình bày như thật và không mang tính châm biếm, được xem là tin giả. Trong một số trường hợp, bài xã luận, phóng sự và bài phanh phui có thể cố ý lan truyền thông tin nhằm lừa dối vì lợi ích tài chính hoặc chính trị {[}4{]}.

Ở cấp độ khái niệm, có thể phân biệt ba loại tin giả {[}3{]}: bịa đặt nghiêm trọng (serious fabrications), tin lừa bịp (hoaxes) và châm biếm (satire). Bịa đặt nghiêm trọng là các bản tin viết dựa trên thông tin sai, kể cả tin đồn về người nổi tiếng. Tin lừa bịp là thông tin sai đưa qua mạng xã hội với mục đích được các nền tảng tin tức truyền thống đăng lại. Châm biếm là dùng hài hước trong tin tức để bắt chước tin thật, nhưng bằng sự mỉa mai và phi lý. Một số nền tảng tin châm biếm nổi tiếng hiện nay là The Onion và The Beaverton, đối lập với các nhà xuất bản tin thật như The New York Times {[}3{]}.

Dù có nhiều nghiên cứu và công cụ phát hiện tin giả, rất ít công trình học thuật tập trung vào nhu cầu thúc đẩy năng lực thông tin (information literacy) để mọi người có thể đánh giá có phản biện thông tin được cung cấp và ra quyết định sáng suốt hơn {[}9{]}.

Stein-Smith {[}5{]} nhấn mạnh rằng năng lực thông tin/truyền thông đã trở thành kỹ năng quan trọng hơn kể từ khi khái niệm tin giả đi vào thảo luận công chúng. Năng lực thông tin không còn là kỹ năng "có thì tốt" mà là yêu cầu để hiểu các tiêu đề tin tức và tham gia thảo luận công cộng. Các cơ sở giáo dục đại học cần mở các khóa học về năng lực thông tin để trang bị cho sinh viên và nhân viên những công cụ cần thiết nhằm nhận diện, chọn lọc, hiểu và sử dụng thông tin đáng tin cậy {[}1{]}. Ngoài môi trường học thuật, năng lực thông tin còn là kỹ năng suốt đời với nhiều ứng dụng trong đời sống hằng ngày {[}5{]}. Các lựa chọn và quan điểm của mỗi người cần dựa trên việc diễn giải đúng thông tin chính xác, kịp thời và có ý nghĩa {[}5{]}.

\subsection{2.3 Tư duy phản biện}

Tư duy phản biện (critical thinking) bao gồm nhiều kỹ năng: suy luận bằng lời, phân tích lập luận, tư duy như kiểm định giả thuyết, xử lý khả năng xảy ra và tính bất định, cùng kỹ năng ra quyết định và giải quyết vấn đề {[}10{]}. Trong nghiên cứu này, vì quan tâm đến việc đánh giá độ tin cậy của tin tức trực tuyến, định nghĩa sau được sử dụng: tư duy phản biện là "khả năng phân tích và đánh giá các lập luận dựa trên tính vững chắc và độ tin cậy của chúng, phản hồi lập luận và rút ra kết luận bằng suy diễn từ thông tin đã cho" {[}11{]}. Nghiên cứu này muốn xem xét cách các kỹ năng nêu trong {[}11{]} có thể được dùng như một phần của năng lực thông tin để nhận diện tin giả tốt hơn.

Phần tiếp theo trình bày phương pháp nghiên cứu được dùng để thực hiện SLR (tổng quan tài liệu có hệ thống).

\section{3 Phương pháp nghiên cứu}

Mục này nêu câu hỏi nghiên cứu, các từ khóa tìm kiếm áp dụng trên cơ sở dữ liệu và các tiêu chí lọc kết quả. Câu hỏi nghiên cứu của SLR (tổng quan hệ thống tài liệu) là:

\begin{itemize}
\item
  Theo các nghiên cứu trước, tư duy phản biện đóng vai trò gì trong việc nhận diện tin giả?
\end{itemize}

Câu hỏi được xác định dựa trên chủ đề nghiên cứu, nhằm xem các nghiên cứu được chọn có cung cấp hiểu biết về việc dùng tư duy phản biện để đánh giá độ tin cậy của tin tức trực tuyến, đặc biệt là nhận diện tin giả hay không.

\textbf{Phạm vi loại trừ.} Theo gợi ý của {[}2{]}, SLR loại các khái niệm tin giả và thuật ngữ liên quan sau:

\begin{itemize}
\item
  Lỗi đưa tin không cố ý;
\item
  Tin đồn không bắt nguồn từ một bài báo cụ thể;
\item
  Thuyết âm mưu;
\item
  Châm biếm ít có khả năng bị hiểu nhầm là sự thật;
\item
  Phát ngôn sai của chính trị gia; và
\item
  Bản tin thiên lệch hoặc gây hiểu lầm nhưng không hoàn toàn sai.
\end{itemize}

\textbf{Từ khóa tìm kiếm.} Nguồn tài liệu được lấy từ Google Scholar (https://scholar.google.com). Chuỗi tìm kiếm được chạy trên Google Scholar, rồi các bài báo và siêu dữ liệu được nạp vào Mendeley để quản lý tài liệu tham khảo. Chuỗi tìm kiếm là:

(``critical think\emph{'' OR ``critically (NEAR/2) reason}'' OR ``critical (NEAR/2) thought\emph{'' OR ``critical (NEAR/2) judge}'' AND ``fake news'' AND (identify* OR analyse* OR find* OR describe* OR review).

Chuỗi này được xây dựng từ chủ đề nghiên cứu (các từ khóa tạo thành tiêu chí cơ sở), sau đó điều chỉnh theo yêu cầu cú pháp của Google Scholar.

\textbf{Tiêu chí chọn lọc.} Các tiêu chí này đặt quy tắc xác định nguồn, thu hẹp tìm kiếm và tập trung vào một chủ đề cụ thể. Tiêu chí đưa vào và loại ra được trình bày ở Bảng 1.

Bảng 1. Tiêu chí đưa vào và loại ra khi chọn bài

{\small
\begin{longtable}[]{@{}
  >{\raggedright\arraybackslash}p{(\columnwidth - 2\tabcolsep) * \real{0.5000}}
  >{\raggedright\arraybackslash}p{(\columnwidth - 2\tabcolsep) * \real{0.5000}}@{}}
\toprule\noalign{}
\begin{minipage}[b]{\linewidth}\raggedright
Tiêu chí đưa vào
\end{minipage} & \begin{minipage}[b]{\linewidth}\raggedright
Tiêu chí loại ra
\end{minipage} \\
\midrule\noalign{}
\endhead
\bottomrule\noalign{}
\endlastfoot
Ấn phẩm liên quan đến "sự thật thay thế", tin giả và thông tin bịa đặt & Ấn phẩm liên quan đến "sự thật thay thế", tin giả và thông tin bịa đặt \\
Tạp chí học thuật về công nghệ thông tin và các lĩnh vực liên quan & Tạp chí học thuật về công nghệ thông tin và các lĩnh vực liên quan \\
Tạp chí học thuật phải nêu tư duy phản biện, kỹ thuật nhận diện tin giả hoặc xem xét tin giả bằng tư duy phản biện & Tạp chí học thuật phải nêu tư duy phản biện, kỹ thuật nhận diện tin giả hoặc xem xét tin giả bằng tư duy phản biện \\
Tạp chí học thuật phải có tóm tắt (abstract) & Tạp chí học thuật phải có tóm tắt (abstract) \\
& Ấn phẩm liên quan đến "sự thật thay thế", tin giả và thông tin bịa đặt \\
\end{longtable}
}

\textbf{Chọn nguồn.} Áp dụng tiêu chí tìm kiếm trên cơ sở dữ liệu trực tuyến thu được 91 bài; các tiêu chí ở Bảng 1 được dùng để thu hẹp còn các bài phù hợp.

\textbf{Lưu đồ PRISMA.} Quá trình chọn lọc gồm bốn giai đoạn (Hình 1). Giai đoạn nhận diện: Google Scholar trả về 91 kết quả, thêm 3 nguồn suy ra từ các nguồn đã xác định, tổng cộng 94 nguồn. Giai đoạn sàng lọc: không có bản trùng; sau khi sàng kỹ, gồm cả xét khả năng tiếp cận bài (dùng miễn phí), còn 39 bản ghi, loại 55 bài. Giai đoạn đủ điều kiện: trong 39 bài, 9 bài bị loại do tiêu đề và tóm tắt không liên quan chủ đề. Danh sách cuối gồm 22 bài đưa vào SLR. Để chuẩn bị phân tích dữ liệu, một bảng trích xuất dữ liệu được lập, phân loại mỗi bài theo: tác giả; tiêu đề; chủ đề (tóm tắt ngắn); năm; quốc gia; và loại ấn phẩm. Bảng này hỗ trợ phân tích các phát hiện ở mục sau.

Hình 1. Lưu đồ PRISMA

\section{4 Phân tích các phát hiện}

\subsection{4.1 Thống kê mô tả}

Do số nghiên cứu liên quan còn ít, việc tìm kiếm không đặt ngày bắt đầu; các bài được đưa vào đến hết 31/8/2019. Phần lớn các bài xuất bản năm 2017 (8 bài) và 2018 (9 bài), phù hợp với việc "fake news" được chọn là từ của năm 2016 {[}8{]}.

Hình 2. Biểu đồ Venn thể hiện mức chồng lấn của các bài theo trọng tâm chính

Các bài báo được chọn được phân loại theo chủ đề. Hình 2 là biểu đồ Venn thể hiện mức độ chồng lấn của các bài báo theo chủ đề trong toàn bộ tổng quan. Nhóm chủ đề "tin giả" có số lần xuất hiện cao nhất, tổng cộng 11 bài. Ba bài tập trung chủ yếu vào "tư duy phản biện", và "năng lực thông tin" là trọng tâm chính của bốn bài. Hai bài kết hợp cả ba chủ đề: tư duy phản biện, năng lực thông tin và tin giả.

Phân tích số bài theo quốc gia cho thấy Hoa Kỳ chiếm ưu thế rõ rệt với 17 bài, tương đương 74\% số bài được chọn. Các quốc gia còn lại có bài đăng gồm Úc, Đức, Ireland, Lebanon, Ả Rập Xê Út và Thụy Điển, mỗi nước một bài.

Xét theo loại ấn phẩm, 15 bài là bài báo tạp chí khoa học, bốn bài là báo cáo, một bài là luận văn, một bài là bài trên tạp chí phổ thông và một bài là trang web.

\textbf{4.2 Thảo luận về các chủ đề}

Các nội dung sau xuất hiện từ phân tích chủ đề (thematic analysis) của các bài báo.

\textbf{Tin giả}

\textbf{Tin giả và trách nhiệm giải trình.} Với ảnh hưởng của mạng xã hội đối với sự lan truyền tin giả {[}2{]}, ai phải chịu trách nhiệm về việc phát tán và tiếp nhận tin giả trong công chúng? Giả định đầu tiên là trong thời đại số, các nền tảng mạng xã hội như Facebook và Twitter nên biên tập thông tin hoặc kiểm chứng sự thật khi bài đăng được tải lên {[}12{]}, nhưng điều này gần với việc xâm phạm quyền tự do ngôn luận. Dù các tác giả đều đồng ý cần có biện pháp hạn chế sự lan truyền của tin giả {[}12, 13{]}, họ khác nhau về việc trách nhiệm thuộc về ai. Metaxas và Mustafaraj {[}13{]} hướng tới phát triển thuật toán hoặc plug-in hỗ trợ xây dựng niềm tin và cho rằng người dùng phải có khả năng nhận diện thông tin sai lệch để tự quyết định có chia sẻ thông tin đó hay không. Ngược lại, Lazer và cộng sự {[}12{]} cho rằng trách nhiệm thuộc về chủ sở hữu nền tảng, phải đặt ra giới hạn đối với loại dữ liệu được phát tán. Do công trình của Metaxas và Mustafaraj {[}13{]} đã thực hiện cách đây bảy năm, có thể kết luận rằng việc dùng thuật toán/plug-in kiểm chứng sự thật chưa thành công trong việc kiềm chế đà lan truyền của tin giả.

\textbf{Tin giả và nghiên cứu của sinh viên.} Có tổng cộng bốn bài tập trung vào nghiên cứu của sinh viên liên quan đến tin giả. Harris, Paskin và Stein-Smith {[}4, 5, 14{]} đều đồng ý rằng sinh viên không có khả năng phân biệt tin thật và tin giả. Một nghiên cứu của Stanford History Education Group cho thấy sinh viên chưa được trang bị để phân biệt tin thật với tin giả {[}4{]}. Phần lớn sinh viên có thể tìm kiếm thông tin đơn giản bằng Google, nhưng không xác định được tác giả của nguồn trực tuyến, hay thông tin đó có gây hiểu lầm hay không {[}14{]}. Hơn nữa, sinh viên không nhận thức được lợi ích của việc học năng lực thông tin ở trường trong việc trang bị kỹ năng cần thiết để nhận diện chính xác tin giả {[}5{]}. Tại Metropolitan Campus của Đại học Fairleigh Dickson, các thủ thư đã đảm nhận vai trò đào tạo kỹ năng năng lực thông tin để nhận diện tin giả {[}5{]}.

\textbf{Tin giả và mạng xã hội.} Một số tác giả {[}6, 15{]} đồng thuận rằng mạng xã hội, nguồn tin hàng đầu, là động lực lớn nhất của tin giả. Mạng xã hội mang lại lợi thế đáng kể để phát tán thông tin bị thao túng; đây là nền tảng mở, không có biên tập viên kiểm duyệt và ai cũng có thể đóng góp nội dung. Theo Nielsen và Graves cũng như Janetzko {[}6, 15{]}, mọi người không nhận diện đúng tin giả; họ có xu hướng gắn tin giả với báo chí kém chất lượng hơn là với thông tin sai được thiết kế để đánh lừa. Hai bài, {[}15{]} và {[}6{]}, bàn về vai trò của tư duy phản biện khi tương tác trên mạng xã hội. Mạng xã hội đưa đến cho ta thông tin đã được lọc theo những gì ta vốn tiêu thụ, khiến người dùng khó tư duy phản biện. Nghiên cứu của Nielsen và Graves {[}6{]} xác nhận rằng việc sinh viên không kiểm chứng được các nguồn trực tuyến sai cần được quan tâm khẩn cấp, vì điều này có thể cho thấy sinh viên là mục tiêu dễ dàng của thông tin bị thao túng.

\textbf{Tin giả chi phối chính trị.} Hai nghiên cứu đề cập tác động của mạng xã hội và sự lan truyền tin giả, cũng như việc chúng có thể đã giúp Donald Trump thắng cử tổng thống Hoa Kỳ năm 2016 {[}2, 16{]}. Ngoài ra, {[}8{]} và {[}2{]} nhắc đến câu chuyện Giáo hoàng ủng hộ Trump trong chiến dịch tranh cử, được chia sẻ rộng rãi (hơn một triệu lần) trên Facebook năm 2016. Các bài này cũng chỉ ra rằng trong thời đại thông tin, việc kiểm chứng sự thật trở nên tương đối dễ, nhưng mọi người có xu hướng tin vào trực giác về các tin tức họ đọc hơn là kiểm tra độ tin cậy của câu chuyện. Việc dùng troll được trả tiền và bot của Nga để tràn ngập nguồn cấp mạng xã hội bằng thông tin sai lệch nhằm xoay chuyển cuộc bầu cử tổng thống Hoa Kỳ theo hướng có lợi cho Donald Trump cũng được nêu rõ {[}16{]}. Việc tạo tin giả với các tiêu đề giật gân ("click bait") tạo ra lượng truy cập khổng lồ vào các trang web gốc, qua đó làm tăng doanh thu quảng cáo {[}2{]}. Điều này có nghĩa là người tạo nội dung bị thúc ép tạo tin giả để tăng doanh thu quảng cáo cho trang web của mình, dù bản thân họ có thể không tin vào tin giả đó {[}2{]}.

\textbf{Năng lực thông tin.} Năng lực thông tin là khi một người có quyền truy cập thông tin, từ đó có thể xử lý những phần mình cần và tạo ra cách sử dụng thông tin tốt nhất {[}1{]}. Dạy sinh viên tầm quan trọng của kỹ năng năng lực thông tin là then chốt, không chỉ để nhận diện tin giả mà còn để định hướng các khía cạnh của cuộc sống đòi hỏi

việc quản lý và thẩm định thông tin, như được bàn trong {[}1, 17{]} và {[}9{]}. Courtney {[}17{]} nhấn mạnh rằng sinh viên báo chí, hơn sinh viên các ngành khác, có thể cần được lồng ghép năng lực thông tin vào chương trình học để nâng cao nhận thức về các câu chuyện tin giả, qua đó hình thành vai trò của người làm tin khách quan và đáng tin cậy. Courtney khảo sát nhiều trường đại học đào tạo báo chí và truyền thông, và xác định rằng sinh viên nhìn chung thiếu nhận thức về mức độ hữu ích của các dịch vụ thư viện trong việc hỗ trợ năng lực thông tin. Courtney {[}17{]} và Rose-Wiles {[}9{]} cho rằng việc sử dụng tài nguyên thư viện cần được bình thường hóa đối với sinh viên. Vì thế hệ millennials và thế hệ Z coi mạng xã hội là điểm tiếp xúc đầu tiên, Rose-Wiles {[}9{]} kêu gọi các trường đại học, cao đẳng và viện nghiên cứu khuyến khích dùng tài nguyên thư viện nhiều hơn tài nguyên trên internet, để sinh viên dựa vào các nguồn đáng tin cậy. Nhìn chung điều này có thể khó thực hiện, nên Rose-Wiles {[}9{]} đề xuất dạy năng lực thông tin và tư duy phản biện để sinh viên vận dụng các kỹ năng này trong mọi tình huống hoặc nguồn thông tin.

Hiện tượng được gọi là "suy thoái sự thật" (truth decay) khiến con người đến mức không còn cần đồng ý với các sự kiện {[}18{]}. Do phân cực chính trị, công chúng cho rằng mình thuộc một nhóm người bị áp bức, nên sẽ tin một nhà lãnh đạo chính trị đánh vào câu chuyện đó {[}18{]}. Cần có hành động cụ thể nhằm thúc đẩy sự tham gia của công dân, khuyến khích mọi người suy nghĩ phản biện, phân tích thông tin và không tin mọi điều mình đọc.

Tư duy phản biện. Chỉ ba bài báo lấy tư duy phản biện làm chủ đề chính. Bronstein và cộng sự {[}19{]} cho rằng một số niềm tin giáo điều và tôn giáo khiến cá nhân có xu hướng tin mọi thông tin được đưa ra mà không cần xem xét kỹ thêm rồi mới quyết định tính xác thực. Bài báo nói rõ thêm rằng những người này cũng dễ tin thuyết âm mưu hơn và có xu hướng hợp lý hóa các sự kiện vô lý. Nghiên cứu của Bronstein và cộng sự {[}19{]} kết luận rằng chủ nghĩa giáo điều và chủ nghĩa chính thống tôn giáo có tương quan cao với việc tin tin giả. Nghiên cứu {[}19{]} đề xuất các biện pháp can thiệp nhằm tăng tư duy cởi mở và tư duy phân tích, như một cách giúp hạn chế niềm tin vào tin giả ở người theo tôn giáo. Howlett {[}20{]} mô tả tư duy phản biện là thực hành dựa trên bằng chứng, tức là đưa lý thuyết về các kỹ năng và khái niệm của tư duy phản biện vào ứng dụng hằng ngày. Jackson {[}21{]} giải thích rằng internet cố ý tự hào là nền tảng của "nội dung chưa qua kiểm duyệt", vì người ta có thể không xem lại nội dung đó, nên nó cần thu hút chú ý ngay lúc này chứ không nhất thiết chính xác. Jackson {[}21{]} nói thêm rằng mạng xã hội ảnh hưởng đến tư duy phản biện khi thay đổi cách nhìn về thông tin đã xuất bản, vốn nay bị xem là các dạng truyền thông thông tin cũ. Điều này đặt ra thách thức cho tư duy phản biện, vì phần lớn thông tin trên internet không chỉ thiếu tin cậy mà còn có thể sai. Jackson {[}21{]} cho rằng một trong những nguy cơ lớn nhất đối với tư duy phản biện có thể là việc mọi người có cảm giác quyền lực khi chỉ cần một tìm kiếm web đơn giản là tìm được những gì mình muốn. Con người không còn quan tâm đến việc đánh giá độ tin cậy của thông tin mình nhận và chia sẻ, từ đó dẫn đến sự lan truyền tin giả {[}21{]}.

\section{5 Thảo luận kết quả}

Dữ liệu tổng hợp trong bài tổng quan này cho thấy tin giả được nhìn nhận ra sao, mức độ quan tâm dành cho nó và những hạn chế của con người khi nhận diện tin giả. Kể từ khi nhận thức về tin giả tăng lên vào năm 2016, mức độ quan tâm của giới học thuật đối với chủ đề này cũng tăng, với phần lớn bài báo được công bố trong giai đoạn 2017-2018. Năm mươi phần trăm số bài tập trung vào tin giả, 18\% bàn về năng lực thông tin và chỉ 13\% bàn về tư duy phản biện.

Phần thảo luận theo chủ đề nhóm và tổng hợp các bài báo theo ba chủ đề chính: tin giả, năng lực thông tin và tư duy phản biện.

\begin{itemize}
\item
  \textbf{Tin giả và trách nhiệm giải trình:} đặt câu hỏi ai phải chịu trách nhiệm về việc lan truyền tin giả, mạng xã hội hay người dùng. Các bài báo kết luận rằng thuật toán kiểm chứng thông tin chưa thành công trong việc hạn chế sự lan truyền tin giả. Phần này cũng đề cập đến tin giả trong nghiên cứu của sinh viên: một nghiên cứu của Stanford History Education Group cho thấy sinh viên chưa được rèn luyện tốt về tư duy phản biện và việc phân biệt tin thật với tin giả {[}4{]}.
\item
  \textbf{Tin giả và mạng xã hội:} cho thấy mạng xã hội vừa là nguồn tin hàng đầu vừa là yếu tố góp phần tạo ra tin giả. Điều này gây khó khăn cho người dùng không có khả năng tư duy phản biện về tin tức trực tuyến hoặc thiếu các kỹ năng cơ bản về năng lực thông tin để nhận diện tin giả.
\item
  \textbf{Tin giả chi phối chính trị:} nhấn mạnh vai trò của tin giả trong chính trị, đặc biệt là cuộc bầu cử tổng thống Mỹ năm 2016 và ảnh hưởng của nó đến cử tri {[}22{]}.
\end{itemize}

Các công bố về năng lực thông tin nhấn mạnh sự cần thiết phải giáo dục công chúng biết nhận diện tin giả, cũng như lợi ích của việc coi năng lực thông tin là một kỹ năng sống {[}1, 9, 17{]}. Nghiên cứu cho thấy sinh viên thường hiểu sai về lợi ích tiềm năng của dịch vụ thư viện. Các tác giả đề xuất thư viện đại học cần được biết đến nhiều hơn và tham gia tích cực hơn với vai trò hỗ trợ, trang bị kỹ năng năng lực thông tin.

Các bài báo về tư duy phản biện chỉ ra hai khía cạnh mà sự thiếu tư duy phản biện khiến người đọc không phân biệt được thông tin đúng và sai. Thứ nhất, sự tự tin vào khả năng tìm kiếm thông tin trực tuyến khiến người ta quá tin vào độ chính xác của thông tin đó {[}21{]}. Thứ hai, thái độ giáo điều và chủ nghĩa cực đoan tôn giáo, vốn khiến người ta tin vào một số tin giả nhất định, gắn với việc thiếu tư duy phản biện và thiếu tinh thần hoài nghi {[}21{]}.

Các bài báo về năng lực thông tin và tư duy phản biện đều thống nhất về giá trị của việc thúc đẩy và giảng dạy các kỹ năng này, đặc biệt cho sinh viên đại học, đối tượng thường là chủ thể của các nghiên cứu được thực hiện.

\section{6 Kết luận}

Bài tổng quan đã xác định 22 bài báo, được tổng hợp và dùng làm bằng chứng để xác định vai trò của tư duy phản biện trong việc nhận diện tin giả. Các bài báo được phân loại theo năm công bố, quốc gia công bố, loại hình công bố và chủ đề. Theo thống kê mô tả, tin giả là xu hướng ngày càng tăng trong những năm gần đây, chủ yếu tại Mỹ kể từ cuộc bầu cử tổng thống năm 2016. Nghiên cứu trong phần lớn các bài báo nhằm đánh giá khả năng nhận diện tin giả của sinh viên. Các nghiên cứu nhất quán trong phát hiện rằng đối tượng nghiên cứu thiếu khả năng phân biệt tin thật và tin giả.

Năng lực thông tin nổi lên như một chủ đề mới từ các nghiên cứu, trong đó Rose-Wiles {[}9{]} khuyến nghị các cơ sở học thuật dạy năng lực thông tin và khuyến khích sinh viên tư duy phản biện khi tiếp cận tin tức trực tuyến. Vai trò tiềm năng của thư viện đại học không chỉ trong việc dạy năng lực thông tin mà còn trong việc hỗ trợ sinh viên đánh giá độ tin cậy của thông tin trực tuyến cũng được nhấn mạnh. Cả ba bài báo đề cập trực tiếp đến tư duy phản biện đều thấy đối tượng nghiên cứu thiếu kỹ năng này, và chỉ ra rằng sự thiếu hụt đó có thể liên quan đến việc con người không nhận diện được tin giả.

Bài tổng quan chỉ ra rằng con người nói chung không có khả năng nhận diện tin giả, đồng thời nhấn mạnh tầm quan trọng của năng lực thông tin và tư duy phản biện như những kỹ năng thiết yếu để đánh giá độ tin cậy của thông tin trực tuyến.

Hạn chế của bài tổng quan là phần lớn nghiên cứu dùng sinh viên làm đối tượng tham gia chính, cho thấy cần chuyển trọng tâm học thuật sang công chúng nói chung. Điều này là cấp thiết vì bất kỳ ai sở hữu thiết bị di động đều có thể là người tạo ra hoặc phát tán tin giả.

Đối với nghiên cứu trong tương lai, đề xuất tiếp tục khảo sát giá trị của việc giảng dạy chính thức về năng lực thông tin tại đại học như một biện pháp giúp sinh viên đánh giá độ tin cậy của tin tức trực tuyến. Do số nghiên cứu về vai trò của tư duy phản biện trong nhận diện tin giả còn rất ít, đây cũng là hướng nghiên cứu quan trọng.
